% ID: SEL-D0071-OB-E03
% ESTADO: validada. S071.
\textbf{Parte I. a)} Por independencia, $P(A\cap B)=0{,}40\cdot0{,}60=\boxed{0{,}24}$. Por inclusión y exclusión, $P(A\cup B)=0{,}40+0{,}60-0{,}24=\boxed{0{,}76}$.

\textbf{b)} La independencia da $P(A\mid B)=\boxed{0{,}40}$. Además, $P(B\cap A^c)=0{,}60-0{,}24=0{,}36$, de modo que $P(B\mid A^c)=0{,}36/0{,}60=\boxed{0{,}60}$.

\textbf{Parte II.} $\hat p=20/80=0{,}25$. Los $20$ éxitos y $60$ fracasos permiten usar la aproximación normal. Para el $90\%$ dejamos un $5\%$ en cada cola y usamos $z_{0{,}95}\approx1{,}64485$. Así,
\[
0{,}25\pm1{,}64485\sqrt{\frac{0{,}25\cdot0{,}75}{80}}
\approx\boxed{[0{,}1704,\ 0{,}3296]}.
\]
